\begin{helper}
For an activated set $A$, priority function $\pi$, and vertices $u,v$, the
pair $\{u,v\}$ is contained in the output independent set of the
\noderef{RandomIndependentSetSampling} priority rule if and only if both
$u$ and $v$ are activated and each beats all of its activated neighbors:
\[
\{u,v\}\subseteq
\{z\in A:\pi(w)<\pi(z)\text{ for every }w\in A\cap N_G(z)\}
\]
if and only if
\[
u,v\in A,\qquad
\pi(w)<\pi(u)\text{ for every }w\in A\cap N_G(u),\qquad
\pi(w)<\pi(v)\text{ for every }w\in A\cap N_G(v).
\]
\end{helper}
\begin{proof}
The output set in \noderef{RandomIndependentSetSampling} is defined by a
vertexwise filter.  If $\{u,v\}$ is a subset of that filter, then applying
the subset relation to $u\in\{u,v\}$ gives that $u\in A$ and that every
activated neighbor of $u$ has smaller priority than $u$.  Applying the same
argument to $v\in\{u,v\}$ gives the analogous two assertions for $v$.

Conversely, assume the four displayed assertions.  To prove
$\{u,v\}$ is contained in the output filter, take any
$z\in\{u,v\}$.  Then either $z=u$ or $z=v$.  In the first case the asserted
activation and priority inequalities for $u$ place $z$ in the filter, and in
the second case the asserted activation and priority inequalities for $v$
place $z$ in the filter.  This proves the desired subset equivalence.
\end{proof}
